\subsection{TENSOR ALGEBRA OF A DIRECT SUM. TENSOR ALGEBRA OF A FREE MODULE. TENSOR ALGEBRA OF A GRADED MODULE}

Let A be a commutative ring and \(M = \bigoplus_{\lambda \in L} M_{\lambda}\) the direct sum of a family of A-modules. It follows from II, §3, no. 7, Proposition 7, by induction on n, that \(T^{n}(M)\) is the direct sum of the submodules the images of the canonical injections

\[
\mathbf {M} _ {\lambda_ {1}} \otimes \mathbf {M} _ {\lambda_ {2}} \otimes \dots \otimes \mathbf {M} _ {\lambda_ {n}} \rightarrow \mathsf {T} ^ {n} (\mathbf {M}) = \mathbf {M} ^ {\otimes n}
\]

relative to all the sequences \((\lambda_{i})\in\mathbf{L}^{n}\) . Identifying \(M_{\lambda_{1}}\otimes M_{\lambda_{2}}\otimes\cdots\otimes M_{\lambda_{n}}\) with this image, it is seen that \(T(\mathbf{M})\) is the direct sum of all the modules

\[
\mathbf {M} _ {\lambda_ {1}} \otimes \mathbf {M} _ {\lambda_ {2}} \otimes \dots \otimes \mathbf {M} _ {\lambda_ {n}}
\]

where \(n\) runs through \(\mathbf{N}\) and, for each \(n\) , \((\lambda_t)\) runs through \(\mathbf{L}^n\) .

We first deduce the following consequence:

THEOREM 1. Let A be a commutative ring, M a free A-module and \((e_{\lambda})_{\lambda \in L}\) a basis of M. Then the elements \(e_{s} = e_{\lambda_{1}} \otimes e_{\lambda_{2}} \otimes \cdots \otimes e_{\lambda_{n}}\) , where \(s = (\lambda_{1}, \ldots, \lambda_{n})\) runs through the set of all finite sequences of elements of L and \(e_{\varnothing}\) is used to denote the unit element of T(M), form a basis of the A-module T(M).

The elements of this basis are obviously homogeneous and the multiplication table is given by

\[
e _ {s} e _ {t} = e _ {s t}\tag{10}
\]

where \(st\) denotes the sequence of elements of \(\mathbf{L}\) obtained by juxtaposing the sequences \(s\) and \(t\) (I, § 7, no. 2).

It is seen that the basis \((e_{s})\) of T(M), with the multiplicative law (10), is canonically isomorphic to the free monoid of the set L (I, §7, no. 2), the isomorphism being obtained by mapping each word s of this monoid to the element \(e_{s}\) . It follows (§2, no. 7) that the tensor algebra T(M) of a free module M, with a basis of indexing set L, is canonically isomorphic to the free associative algebra of L over A. In particular (§2, no. 7, Proposition 7), for every mapping \(f: L \to E\) from L to an A-algebra E, there exists one and only one A-algebra homomorphism \(\bar{f}: T(M) \to E\) such that \(\bar{f}(e_{\lambda}) = f(\lambda)\) .

Remark. The above results can equally be obtained as a consequence of the universal properties of the free associative algebra and of the tensor algebra, using II, § 3, no. 7, Corollary 2 to Proposition 7.

COROLLARY. If M is a projective A-module, T(M) is a projective A-module.

M is a direct factor of a free A-module N (II, §2, no. 2, Proposition 4) and hence T(M) is a direct factor of T(N) (no. 2); as T(N) is free (Theorem 1), this shows that T(M) is projective (II, §2, no. 2).

PROPOSITION 7. Let \(\Delta\) be a commutative monoid, \(\mathbf{M}\) a graded A-module of type \(\Delta\) and \((\mathbf{M}_{\alpha})_{\alpha \in \Delta}\) its graduation. For every ordered pair \((\alpha, n) \in \Delta \times \mathbf{N}\) , let \(\mathsf{T}^{\alpha,n}(\mathbf{M})\) be the (direct) sum of the submodules \(\mathbf{M}_{\alpha_1} \otimes \mathbf{M}_{\alpha_2} \otimes \dots \otimes \mathbf{M}_{\alpha_n}\) of \(\mathsf{T}^n(\mathbf{M})\) such that \(\sum_{i=1}^{n} \alpha_i = \alpha\) ; then \((T^{\alpha, n}(M))_{(\alpha, n) \in \Delta \times N}\) is the only graduation of type \(\Delta \times N\) compatible with the algebra structure on \(T(M)\) and inducing on \(M = T^1(M)\) the given graduation.

It has been seen at the beginning of this no. that \(\mathsf{T}(\mathbf{M})\) is the direct sum of the \(\mathsf{T}^{\alpha,n}(\mathbf{M})\) and the fact that this is a graduation compatible with the algebra structure follows immediately from the definitions. If \((\mathsf{T}'^{\alpha,n})\) is another graduation of type \(\Delta \times \mathbf{N}\) on \(\mathsf{T}(\mathbf{M})\) compatible with the algebra structure and such that \(\mathsf{T}^{\alpha,1}(\mathbf{M}) = \mathsf{T}'^{\alpha,1}\) for \(\alpha \in \Delta\) , it follows immediately from the definitions that, for all \(n \geqslant 1\) and all \(\alpha \in \Delta\) , \(\mathsf{T}^{\alpha,n}(\mathbf{M}) \subset \mathsf{T}'^{\alpha,n}\) ; but since \(\mathsf{T}(\mathbf{M})\) is also the direct sum of the \(\mathsf{T}^{\alpha,n}(\mathbf{M})\) , this implies that \(\mathsf{T}'^{\alpha,n} = \mathsf{T}^{\alpha,n}(\mathbf{M})\) (II, § 1, no. 8, Remark).

